\chapter{Metodologi Penelitian}

\section{Alur Penelitian}
\TemplateTodo{Jelaskan tahapan penelitian dari awal hingga akhir. Pastikan alur sesuai dengan rumusan masalah dan hasil yang akan dilaporkan.}

Metodologi harus cukup rinci agar penelitian dapat dipahami dan, sejauh mungkin, direplikasi. Tuliskan keputusan penting, parameter, asumsi, serta alasan pemilihan metode.

\section{Objek dan Data Penelitian}
\TemplateTodo{Jelaskan objek penelitian, sumber data, kriteria inklusi/eksklusi, periode pengambilan data, jumlah data, dan aspek etik atau perizinan jika relevan.}

\section{Perancangan Metode atau Sistem}
\TemplateTodo{Uraikan rancangan metode, model, sistem, atau prosedur eksperimen. Hubungkan setiap komponen dengan tujuan penelitian.}

\begin{algorithm}[H]
\caption{Contoh format algoritma penelitian}
\KwInput{Data atau masukan penelitian}
\KwOutput{Hasil yang dievaluasi}
Lakukan tahap persiapan data\;
Terapkan metode yang diusulkan\;
Evaluasi keluaran menggunakan metrik yang sesuai\;
\KwReturn{Ringkasan hasil dan analisis}
\end{algorithm}

\section{Skenario Pengujian}
\TemplateTodo{Jelaskan rancangan eksperimen, pembagian data, pembanding, metrik evaluasi, perangkat yang digunakan, dan cara memastikan hasil adil.}

\section{Teknik Analisis Data}
\TemplateTodo{Jelaskan bagaimana hasil akan dianalisis: statistik deskriptif, perbandingan metrik, analisis kesalahan, validasi pengguna, atau pendekatan lain yang sesuai.}
